\chapter{Explicação da tomada de decisão}
\label{cap:explicacao}

Uma das características do sistema pericial é a possibilidade de compreender
o raciocínio que conduziu a uma determinada conclusão e, por isso, permitir
relacionar a decisão com os factos e as regras utilizados durante o processo
de inferência. No contexto deste projeto, a ação selecionada é importante;
porém, interessa perceber também quais as condições que levaram a essa
escolha e por que motivo outras alternativas não prevaleceram. Para
estruturar esta componente, foram consideradas três formas de explicação
complementares: \emph{Why}, \emph{How} e \emph{Why Not}, cada uma orientada
para uma questão diferente relacionada com a tomada da decisão, como
representado na \cref{fig:explicacao}.

\begin{figure}[H]
  \centering
  \begin{tikzpicture}[
      caixa/.style={draw=corEscura, fill=corFundo, thick, rounded corners=2pt,
        text width=9.2cm, minimum height=1.15cm, align=center, font=\small},
      seta/.style={-latex, thick, draw=corEscura}]
    \node[draw=corEscura, fill=corFundo, thick, rounded corners=2pt,
      minimum height=1.15cm, text width=2.4cm, align=center, font=\small]
      (decisao) at (0,0) {Decisão final};
    \node[caixa] (why) at (7.4,1.6)
      {\emph{Why?} -- Porque foi escolhida?\\Regra + evidência};
    \node[caixa] (how) at (7.4,0)
      {\emph{How?} -- Como foi obtida?\\Sensor $\rightarrow$ evento
       $\rightarrow$ facto $\rightarrow$ regra $\rightarrow$ decisão};
    \node[caixa] (whynot) at (7.4,-1.6)
      {\emph{Why Not?} -- Porque não outra ação?\\Condições não satisfeitas /
       alternativa de menor prioridade};
    \draw[seta] (decisao.east) -- (why.west);
    \draw[seta] (decisao.east) -- (how.west);
    \draw[seta] (decisao.east) -- (whynot.west);
  \end{tikzpicture}
  \caption{Três dimensões de explicação}
  \label{fig:explicacao}
\end{figure}

\section{\emph{Why}}
\label{sec:why}

A explicação do \emph{Why} procura dar resposta à questão: \enquote{O porquê de ter
sido tomada esta decisão?}. Neste sentido, é necessário identificar a regra
que originou a proposta, os factos que satisfizeram as respetivas condições e
a prioridade associada à ação. Por exemplo, caso o sistema determine uma
travagem por ter detetado um peão, a explicação deverá permitir perceber que
a decisão foi o resultado da presença do peão no caminho, da distância
observada, do nível de risco identificado e da regra correspondente a essa
situação. Como proposta, passível de alteração, procurar-se-á utilizar uma
estrutura semelhante à seguinte:

\begin{tcolorbox}[enhanced, colback=corFundo, colframe=corEscura,
    boxrule=0pt, leftrule=1.2pt, arc=0pt, outer arc=0pt,
    left=8pt, right=8pt, top=4pt, bottom=4pt, fontupper=\small]
  \textbf{Decisão:} Travar.\\
  \textbf{Motivo:} Foi detetado um peão na trajetória do veículo e a situação
  identificada levou à aplicação de uma regra de segurança cuja proposta
  prevaleceu sobre as restantes.
\end{tcolorbox}

Na implementação atual, as decisões propostas permitem associar informação
relativa à razão apresentada, à regra de origem e à evidência utilizada.
Desta forma, o \emph{Why} permite justificar a ação selecionada através do
conhecimento efetivamente utilizado durante a tomada de decisão.

\section{\emph{How}}
\label{sec:how}

Enquanto o \emph{Why} tem como finalidade justificar a decisão final, o
\emph{How} tem como foco descrever como é que o sistema chegou até essa
conclusão. Neste caso, o objetivo é reconstruir o percurso desde a informação
recolhida até à ação selecionada. A explicação pode, por exemplo, identificar
a origem da informação, o facto criado a partir dessa observação, a regra que
utilizou esse facto e a decisão resultante desse processo. A título de
exemplo, num cenário em que é detetado um peão, a explicação poderá mostrar
que a informação foi obtida através de um sensor, posteriormente representada
como um facto na memória de trabalho, utilizada por uma determinada regra e,
por fim, transformada numa proposta de ação. A possibilidade de consultar
alterações na memória de trabalho e acompanhar as regras que foram ativadas
permite reconstruir este percurso e, consequentemente, perceber de que forma
a conclusão foi obtida.

\section{\emph{Why Not}}
\label{sec:whynot}

O \emph{Why Not} procura dar resposta à seguinte questão: \enquote{Porque não foi
escolhida outra alternativa?}.

Uma alternativa pode não ter sido escolhida por dois motivos principais. Em
primeiro lugar, não reuniu as condições necessárias para ativar a regra que a
escolhesse. Por outro lado, essa ação pode ter sido efetivamente proposta,
mas ter sido recusada face a outra decisão com uma prioridade superior. A
título de exemplo, imagine-se que, perante uma situação de risco, o sistema
poderá considerar uma redução de velocidade e, simultaneamente, uma travagem.
Caso a travagem esteja associada a um nível de prioridade superior, a
explicação deverá indicar que a redução de velocidade não foi selecionada
porque a situação exigia uma ação com maior criticidade. Uma explicação deste
tipo poderá ser feita da seguinte forma:

\begin{tcolorbox}[enhanced, colback=corFundo, colframe=corEscura,
    boxrule=0pt, leftrule=1.2pt, arc=0pt, outer arc=0pt,
    left=8pt, right=8pt, top=4pt, bottom=4pt, fontupper=\small]
  \textbf{Alternativa:} Desacelerar.\\
  \textbf{Motivo de ter sido rejeitada:} Foi identificada uma situação de
  maior risco e respetiva prioridade superior (regra de travagem).
\end{tcolorbox}

O \emph{Why Not} pode também ser utilizado quando uma determinada ação nem
chegou a ser proposta. Por exemplo, se for questionado por que motivo o
veículo não mudou de via, a explicação poderá indicar que a via adjacente se
encontrava ocupada ou que não estavam reunidas as condições necessárias para
realizar a manobra.

No estado atual do projeto, esta componente encontra-se ainda em
desenvolvimento. O sistema consegue justificar a decisão selecionada, mas uma
explicação completa de \emph{Why Not} exige a conservação explícita das
propostas não selecionadas ou das condições que impediram determinadas regras
de produzir uma proposta.

A informação necessária para construir estas explicações poderá incluir a
identificação da regra, os factos e eventos considerados, as condições
verificadas, a prioridade atribuída e as alternativas analisadas. Para além de
facilitar a compreensão das decisões, esta informação poderá também
contribuir para a identificação de regras incompletas ou comportamentos
inesperados durante a validação do sistema.

Desta forma, os mecanismos \emph{Why}, \emph{How} e \emph{Why Not}
complementam-se: o primeiro justifica a decisão tomada, o segundo apresenta o
percurso que conduziu até essa decisão e o terceiro tenta explicar por que
razão uma alternativa não foi selecionada. Em conjunto, estes propósitos
contribuem para um processo de inferência mais transparente e facilitam a
análise das decisões tomadas.
